% =====================================================================
%  PAPER 2  —  A Ruppeiner Curvature Dichotomy for Generalized Prime Gases
%  IG-PRIMON-T1.  Skeleton v0.1 (2026-06-14).
%  Status convention in margins: [V] verified, [E] extrapolation, [C] conjecture.
%  Locked math is filled in; prose marked  %TODO  is to be written.
%  For J. Phys. A submission, swap \documentclass to {iopart} and adjust.
% =====================================================================
\documentclass[a4paper,11pt]{article}
\usepackage{amsmath,amssymb,amsthm,mathtools}
\usepackage{booktabs}
\usepackage[margin=1in]{geometry}
\usepackage[colorlinks=true,linkcolor=blue,citecolor=blue]{hyperref}

\theoremstyle{plain}
\newtheorem{theorem}{Theorem}
\newtheorem{lemma}{Lemma}
\newtheorem{proposition}{Proposition}
\theoremstyle{definition}
\newtheorem{definition}{Definition}
\theoremstyle{remark}
\newtheorem{remark}{Remark}

\newcommand{\Ric}{R}                 % scalar curvature
\newcommand{\Li}{\operatorname{Li}}
\newcommand{\Hess}{\operatorname{Hess}}
\newcommand{\eps}{\varepsilon}
\newcommand{\detg}{\det g}

\title{A Ruppeiner curvature dichotomy for generalized prime gases:\\
       asymptotic flatness at limiting-temperature transitions}
\author{Aaron Jones\\ \small JtechAI LLC}
\date{\today}

\begin{document}
\maketitle

\begin{abstract}
%TODO prose. Target content:
We study the thermodynamic (Ruppeiner) geometry of grand-canonical
\emph{generalized prime gases}, whose partition functions are Dirichlet series
$\log\Xi(\beta,z)=\sum_{k\ge1}(z^k/k)\,\mathcal P(k\beta)$. We prove that the
scalar curvature at a phase transition is governed by \emph{where the singularity
lives}: a \emph{temperature-driven} (Hagedorn / limiting-temperature) transition
forces the curvature to vanish asymptotically, $\Ric\to0$, while a
\emph{fugacity-driven} (condensation) transition makes it diverge, $|\Ric|\to\infty$.
The discriminating mechanism is a single algebraic condition on the singular sector
of $\log\Xi$ --- it is curvature-flat precisely when that sector factorizes as
$e^{y}\varphi(x)$ in the natural coordinates, i.e.\ when fugacity and temperature do
not couple. The result is established by an explicit cancellation lemma together with
an analyticity argument, and is bracketed by four exactly-computable instances: the
classical ideal gas ($\Ric=0$), the Riemann primon gas and an all-integer Bose gas
(both Hagedorn, $\Ric\to0$), and the ideal Bose gas at condensation
($|\Ric|\to\infty$). Two of the four contain no primes; the statement is one of
statistical mechanics.
\end{abstract}

% =====================================================================
\section{Introduction}\label{sec:intro}
%TODO prose. Beats to hit:
% - Ruppeiner geometry: metric = -Hessian of entropy; |R| ~ correlation volume;
%   R -> infinity at ordinary critical points; R = 0 for the classical ideal gas
%   (no interactions). Cite Ruppeiner RMP [1]; Janyszek-Mrugala [2].
% - The primon / Riemann gas: a gas whose modes are the primes, energies log p,
%   partition function zeta(beta), with a Hagedorn (limiting) temperature at beta=1.
%   Cite Julia [3], Spector [4], Bost-Connes [5], Hagedorn [6].
% - The question: a Hagedorn transition has NO diverging correlation length (it is a
%   density-of-states wall, not fluctuation-driven). What does Ruppeiner curvature do
%   there? Folklore vs. fact.
% - This paper's answer: a clean dichotomy with an algebraic discriminant. State it.
% - Emphasize: no number theory enters the theorem; the primon gas is one example.

% =====================================================================
\section{Setup}\label{sec:setup}

\begin{definition}[generalized prime gas]\label{def:gas}
%TODO prose framing. The object:
A grand-canonical gas with
\begin{equation}\label{eq:logXi}
   \log\Xi(\beta,z)=\sum_{k\ge1}\frac{z^k}{k}\,\mathcal P(k\beta),
\end{equation}
where $\mathcal P$ is the spectral (prime) zeta of the mode spectrum, $\beta$ the
inverse temperature and $z$ the fugacity.
\end{definition}

Natural coordinates $(x,y)=(-\beta,\log z)$. The family is exponential with sufficient
statistics $(E,N)$, so the Ruppeiner/Fisher metric is the Hessian of the potential
$\psi:=\log\Xi$,
\begin{equation}\label{eq:metric}
   g=\Hess_{(x,y)}\psi,\qquad
   \partial_x^{a}\partial_y^{b}\psi=(-1)^{a}\sum_{k\ge1}k^{\,a+b-1}z^{k}\,\mathcal P^{(a)}(k\beta).
\end{equation}
For a two-dimensional Hessian metric the scalar curvature is
\begin{equation}\label{eq:R}
   \Ric=-\frac{N}{2(\detg)^{2}},\qquad
   N=\det\!\begin{pmatrix}
       \psi_{xx} & \psi_{xy} & \psi_{yy}\\
       \psi_{xxx}& \psi_{xxy}& \psi_{xyy}\\
       \psi_{xxy}& \psi_{xyy}& \psi_{yyy}
     \end{pmatrix}.
\end{equation}
%TODO remark: formula+sign certified against the normal-family closed form (R=-1) and a
%flat product control (R=0); sphere-positive convention. (Companion: moduleC_certify.py.)

\paragraph{Hypotheses.}
\begin{itemize}
\item[\textbf{(H1)}] \emph{Temperature-driven singularity.} $\mathcal P(s)$ is real-analytic
on $(1,\infty)$ with an isolated singularity at its abscissa $s=1$ and no other singularity
on the real axis $s\ge1$. (The transition is a singularity in $\beta$ at the limiting
temperature $\beta=1$.)
\item[\textbf{(H2)}] \emph{Finite background.} $\sum_{k\ge2}k(k-1)\mathcal P^{(j)}(k)<\infty$
for $j=0,1,2$ (automatic when $\mathcal P(k)\to0$ geometrically).
\end{itemize}

% =====================================================================
\section{The product-form criterion}\label{sec:lemmas}

\begin{lemma}[product form annihilates the curvature numerator]\label{lem:C1}
If $\psi_{\mathrm{sing}}=e^{y}\varphi(x)$ for a smooth $\varphi$, then the third row of $N$
in \eqref{eq:R}, $(\psi_{xxy},\psi_{xyy},\psi_{yyy})$, equals the first row
$(\psi_{xx},\psi_{xy},\psi_{yy})$ identically and at all orders; hence the contribution of
$\psi_{\mathrm{sing}}$ to $N$ is $\equiv0$.
\end{lemma}
\begin{proof}
$\partial_y\!\big(e^{y}\varphi^{(j)}(x)\big)=e^{y}\varphi^{(j)}(x)$, so each entry of row three
equals the corresponding entry of row one. A determinant with two equal rows vanishes. The
identity is exact in $\eps$, so it holds order by order. %TODO: one line on "all orders".
\end{proof}

\begin{lemma}[only the single-particle term is singular]\label{lem:k1}
Under \textnormal{(H1)}, every $k\ge2$ term of \eqref{eq:logXi} is real-analytic at
$(\beta,z)=(1,1)$, and the unique singular term is
\[
   \psi_1=z\,\mathcal P(\beta)=e^{y}\varphi(x),\qquad \varphi(x):=\mathcal P(-x).
\]
\end{lemma}
\begin{proof}
For $k\ge2$ and $\beta$ near $1$, $\mathcal P(k\beta)$ is evaluated at $k\beta\approx k\ge2$,
inside the analytic region by (H1); the term is analytic. Only $k=1$ probes the abscissa,
and it is linear in $z$. The product form is forced by the singularity living in temperature
and being sampled at the integer multiples $k\beta$. %TODO: prose on why this is generic.
\end{proof}

% =====================================================================
\section{The dichotomy theorem}\label{sec:thm}

\begin{theorem}[temperature-driven $\Rightarrow$ asymptotically flat]\label{thm:flat}
Under \textnormal{(H1)--(H2)}, as $(\beta,z)\to(1,1)$ along $z=1$,
\[
   \Ric\longrightarrow 0 .
\]
Explicitly, with $L=\log(1/\eps)$, $\eps=\beta-1$, and background amplitude
$\Delta_3:=\sum_{k\ge2}k(k-1)z^{k}\mathcal P(k)$:
\[
   \mathcal P(1+\eps)=\log(1/\eps)+O(1)\ \Rightarrow\ \Ric\sim\frac{\Delta_3}{2z^{2}(L+\kappa)^{2}},
   \qquad
   \mathcal P(1+\eps)=c/\eps+O(1)\ \Rightarrow\ \Ric\sim \frac{\Delta_3\,\eps^{2}}{z^{2}} .
\]
\end{theorem}
\begin{proof}
By Lemma~\ref{lem:k1} the singular sector is $\psi_1=e^{y}\varphi(x)$; by Lemma~\ref{lem:C1}
it contributes $0$ to $N$. Thus $N$ is sourced entirely by the analytic $k\ge2$ background.
Row-reduce $N$ by $\text{Row}_3\to\text{Row}_3-\text{Row}_1=(\Delta_1,\Delta_2,\Delta_3)+O(\eps)$
(finite by (H2)). The metric determinant is dominated by the singular sector and diverges; the
singular powers of $\eps$ cancel between $N$ and $(\detg)^2$, leaving a residual factor that
tends to zero --- $1/(L+\kappa)^2$ in the log case (full computation in
Appendix~\ref{app:rates}), $\eps^2$ in the simple-pole case. %TODO: fold the explicit
% cofactor computation; cite the C-derivation for the log case.
\end{proof}

\begin{proposition}[fugacity-driven $\Rightarrow$ divergent]\label{prop:div}
If the transition is at $z\to z_c$ with the singularity carried by the full fugacity series,
$\log\Xi_{\mathrm{sing}}=\varphi(\beta)\,\Li_a(z)$, then the singular sector
$\varphi(\beta)\,\Li_a(e^{y})$ is not of the form $e^{y}\varphi(x)$; Lemma~\ref{lem:C1} does
not apply, $\text{Row}_3\ne\text{Row}_1$, and $|\Ric|\to\infty$.
\end{proposition}
%TODO proof/sketch: the standard |R| ~ correlation-volume divergence, recovered as the complement.

\begin{remark}\label{rem:criterion}
Theorem~\ref{thm:flat} and Proposition~\ref{prop:div} together give the discriminant in one line:
\[
   \boxed{\ \Ric\to0 \iff \text{the singular sector of }\log\Xi\text{ factorizes as } e^{y}\varphi(x).\ }
\]
%TODO prose: physical reading -- product form = a single dominant mode drives the transition,
% fugacity factors out as a bare z, no diverging correlation length; broken product form =
% fluctuation-driven criticality with field/temperature coupled through a scaling function.
\end{remark}

% =====================================================================
\section{Instances}\label{sec:instances}
%TODO prose intro: all four run through the same engine \eqref{eq:R}; the classical case
% validates it (R=0 is forced and obtained).

\begin{table}[h]
\centering
\begin{tabular}{@{}llll@{}}
\toprule
gas & singular structure & type & curvature (rate / amplitude)\\
\midrule
classical ideal gas        & $u^{-3/2}e^{y}\ (\varphi\cdot e^{y})$ & product form     & $\Ric=0$ (to $10^{-31}$)\\
Riemann primon gas         & $z\,P(\beta),\ P\sim\log(1/\eps)$    & temp-driven, log & $\Ric\to0^{+}$, $\sim1/(L+\kappa)^2$, $\Delta_3=5.0451882$\\
all-integer Bose gas       & $z(\zeta(\beta)-1)\sim1/\eps$        & temp-driven, pole& $\Ric\to0^{+}$, $\Ric/\eps^2\to\Delta_3=5.693982$\\
ideal Bose gas (BEC)       & $u^{-3/2}\Li_{5/2}(e^{y})$           & fugacity-driven  & $|\Ric|\to\infty$, $\sim\tau^{-1/2}$\\
\bottomrule
\end{tabular}
\caption{The dichotomy, bracketed. $u=\beta$; $\tau=-\log z\to0^{+}$ at condensation;
$\kappa=0.832503$ for the primon gas. The two flat instances differ in spectrum (primes vs.\
all integers) and in singularity type (log vs.\ pole); two of the four contain no primes.}
\label{tab:instances}
\end{table}

\paragraph{All-integer Bose gas.} $\prod_{n\ge2}(1-z\,n^{-\beta})^{-1}$, i.e.\ $\mathcal P=\zeta-1$,
along $z=1$:
\[
\begin{array}{r|cccc}
\eps      & 10^{-2}      & 10^{-3}      & 10^{-4}      & \to0\\\hline
\Ric/\eps^2 & 4.95\!\dots & 5.61\!\dots & 5.69\!\dots & \Delta_3^{\mathrm{int}}=5.693982
\end{array}
\]
%TODO: ideal-Bose-gas table (R climbing 0.33->3.59 as tau:0.3->0.001, R*sqrt(tau)->~0.11).

\begin{remark}[a geometric invariant of the primon gas]\label{rem:Cconst}
The renormalized information length of the Riemann primon gas from the Hagedorn point to
$\beta=2$,
\[
   C=\int_{1}^{2}\!\Big[\sqrt{\partial_\beta^2\log\zeta(\beta)}-\tfrac{1}{\beta-1}\Big]\,d\beta
    =-0.0343561541791219860831108814584476\ldots,
\]
is, to $90$ digits, not an integer combination of $\{\pi,\log2,\log3,\gamma,\gamma_1,\zeta(3),
\zeta(5),G,\log A,\dots\}$ at coefficient height $\le10^{6}$ (PSLQ); we register it as a new
constant of the gas. %TODO: decide whether this belongs here or is deferred to Paper 1.
\end{remark}

% =====================================================================
\section{Discussion}\label{sec:disc}
%TODO prose. Beats:
% - The criterion as physics: flatness <-> no diverging correlation length.
% - Scope: theorem for the class \eqref{eq:logXi} under (H1)-(H2); NOT all limiting-temperature
%   transitions in nature. State this plainly.
% - The general-rate gap (arbitrary analytic singularity profile): CLOSED. Appendix A's cofactor
%   reduction derives the general pole-order rate R ~ Delta_3(alpha+1)/(2 A^2) eps^{2 alpha} (and the
%   log case); no longer an open item. [Superseded by paper2_curvature_dichotomy_FULL_v0_1.tex.]
% - No bearing on zeros / RH; deliberately a stat-mech statement.

% =====================================================================
\appendix
\section{Curvature rates: the cofactor computation}\label{app:rates}
%TODO: the explicit Row3-Row1 reduction and cofactor orders giving 1/(L+kappa)^2 (log) and
% eps^2 (pole). Numerics reproduced from the contrast runs.

% =====================================================================
\begin{thebibliography}{9}
%TODO verify exact pagination on all entries.
\bibitem{ruppeiner} G.~Ruppeiner, \emph{Riemannian geometry in thermodynamic fluctuation theory},
   Rev.\ Mod.\ Phys.\ \textbf{67}, 605 (1995).
\bibitem{janyszek} H.~Janyszek and R.~Mrug{\l}a, \emph{Riemannian geometry and stability of ideal
   quantum gases}, J.\ Phys.\ A \textbf{23}, 467 (1990).
\bibitem{julia} B.~L.~Julia, \emph{Statistical theory of numbers}, in \emph{Number Theory and
   Physics}, eds.\ J.-M.~Luck et al.\ (Springer, 1990).
\bibitem{spector} D.~Spector, \emph{Supersymmetry and the M\"obius inversion function},
   Commun.\ Math.\ Phys.\ \textbf{127}, 239 (1990).
\bibitem{bostconnes} J.-B.~Bost and A.~Connes, \emph{Hecke algebras, type III factors and phase
   transitions with spontaneous symmetry breaking in number theory}, Selecta Math.\ \textbf{1}, 411 (1995).
\bibitem{hagedorn} R.~Hagedorn, \emph{Statistical thermodynamics of strong interactions at high
   energies}, Nuovo Cimento Suppl.\ \textbf{3}, 147 (1965).
\end{thebibliography}

\end{document}
